\begin{textsample}{POS Dim 1 – 2011 – Score 23.00 – 2011/t002559.txt}  \label{ex:f1_pos_033}
Each of you possesses \textbf{the} most powerful, dangerous and subversive trait that natural \textbf{selection} has ever \textbf{devised}.
It ' s a piece of neural audio technology for rewiring other people ' s minds.
I ' m talking about your language, of course, because it allows you to implant a thought from your mind directly into someone else ' s mind, and they can attempt to do \textbf{the} same to you, without either of you having to perform surgery.
Instead, when you speak, you ' re actually using a form of telemetry not so different from \textbf{the} remote control \textbf{device} for your television.
It ' s just that, whereas that \textbf{device} relies on \textbf{pulses} of infrared light, your language relies on \textbf{pulses}, discrete \textbf{pulses}, of sound.
And just as you use \textbf{the} remote control \textbf{device} to alter \textbf{the} internal settings of your television to suit your mood, you use your language to alter \textbf{the} settings inside someone else ' s brain to suit your interests.
Languages are genes talking, getting things that they want.
And just imagine \textbf{the} sense of wonder in a baby when it first discovers that, merely by uttering a sound, it can get \textbf{objects} to move across a room as if by magic, and maybe even into its mouth.
Now language ' s subversive power has been recognized throughout \textbf{the} ages in censorship, in books you can ' t read, phrases you can ' t use and words you can ' t say.
In fact, \textbf{the} Tower of Babel story in \textbf{the} Bible is a fable and warning about \textbf{the} power of language.
According to that story, early humans developed \textbf{the} conceit that, by using their language to work together, they could build a tower that would take them all \textbf{the} way to heaven.
Now God, angered at this attempt to usurp his power, destroyed \textbf{the} tower, and then to ensure that it would never be rebuilt, he scattered \textbf{the} people by giving them different languages — confused them by giving them different languages.
And this leads to \textbf{the} wonderful irony that our languages exist to prevent us from communicating.
Even today, we know that there are words we cannot use, phrases we cannot say, because if we do so, we might be accosted, jailed, or even killed.
And all of this from a puff of air emanating from our mouths.
Now all this fuss about a single one of our traits tells us there ' s something worth explaining.
And that is how and why did this remarkable trait \textbf{evolve}, and why did it \textbf{evolve} only in our species?
Now it ' s a little bit of a surprise that to get an answer to that question, we have to go to tool use in \textbf{the} chimpanzees.
Now these chimpanzees are using tools, and we take that as a sign of their intelligence.
But if they really were intelligent, why would they use a stick to extract termites from \textbf{the} ground rather than a shovel?
And if they really were intelligent, why would they crack open \textbf{nuts} with a \textbf{rock}?
Why wouldn ' t they just go to a shop and buy a bag of \textbf{nuts} that somebody else had already cracked open for them?
Why not?
I mean, that ' s what we do.
Now \textbf{the} reason \textbf{the} chimpanzees don ' t do that is that they lack what psychologists and anthropologists call social \textbf{learning}.
They seem to lack \textbf{the} ability to learn from others by \textbf{copying} or imitating or simply watching.
As a result, they can ' t improve on others ' ideas or learn from others ' mistakes — benefit from others ' wisdom.
And so they just do \textbf{the} same thing over and over and over again.
In fact, we could go away for a million years and come back and these chimpanzees would be doing \textbf{the} same thing with \textbf{the} same sticks for \textbf{the} termites and \textbf{the} same \textbf{rocks} to crack open \textbf{the} \textbf{nuts}.
Now this may sound arrogant, or even full of hubris.
How do we know this?
Because this is exactly what our ancestors, \textbf{the} Homo erectus, did.
These upright apes \textbf{evolved} on \textbf{the} African savanna about two million years ago, and they made these splendid hand axes that fit wonderfully into your hands.
But if we look at \textbf{the} \textbf{fossil} record, we see that they made \textbf{the} same hand axe over and over and over again for one million years.
You can follow it through \textbf{the} \textbf{fossil} record.
Now if we make some guesses about how long Homo erectus lived, what their generation time was, that ' s about 40, 000 generations of parents to offspring, and other individuals watching, in which that hand axe didn ' t change.
It ' s not even clear that our very close genetic \textbf{relatives}, \textbf{the} Neanderthals, had social \textbf{learning}.
Sure enough, their tools were more complicated than those of Homo erectus, but they too showed very little change over \textbf{the} 300, 000 years or so that those species, \textbf{the} Neanderthals, lived in Eurasia. \textbf{Okay}, so what this tells us is that, contrary to \textbf{the} old adage,"monkey see, monkey do,"\textbf{the} surprise really is that all of \textbf{the} other animals really cannot do that — at least not very much.
And even this picture has \textbf{the} suspicious taint of being rigged about it — something from a Barnum \& Bailey circus.
But by comparison, we can learn.
We can learn by watching other people and \textbf{copying} or imitating what they can do.
We can then choose, from among a range of options, \textbf{the} best one.
We can benefit from others ' ideas.
We can build on their wisdom.
And as a result, our ideas do accumulate, and our technology progresses.
And this cumulative cultural adaptation, as anthropologists call this accumulation of ideas, is responsible for everything around you in your bustling and teeming everyday lives.
I mean \textbf{the} world has changed out of all proportion to what we would recognize even 1, 000 or 2, 000 years ago.
And all of this because of cumulative cultural adaptation. \textbf{The} chairs you ' re sitting in, \textbf{the} lights in this auditorium, my microphone, \textbf{the} iPads and iPods that you carry around with you — all are a result of cumulative cultural adaptation.
Now to many commentators, cumulative cultural adaptation, or social \textbf{learning}, is job done, end of story.
Our species can make stuff, therefore we prospered in a way that no other species has.
In fact, we can even make \textbf{the}"stuff of life"— as I just said, all \textbf{the} stuff around us.
But in fact, it turns out that some time around 200, 000 years ago, when our species first arose and acquired social \textbf{learning}, that this was really \textbf{the} beginning of our story, not \textbf{the} end of our story.
Because our acquisition of social \textbf{learning} would create a social and \textbf{evolutionary} dilemma, \textbf{the} resolution of which, it ' s fair to say, would determine not only \textbf{the} future course of our psychology, but \textbf{the} future course of \textbf{the} entire world.
And most importantly for this, it ' ll tell us why we have language.
And \textbf{the} reason that dilemma arose is, it turns out, that social \textbf{learning} is visual theft.
If I can learn by watching you, I can steal your best ideas, and I can benefit from your efforts, without having to put in \textbf{the} time and energy that you did into developing them.
If I can watch which lure you use to catch a \textbf{fish}, or I can watch how you flake your hand axe to make it better, or if I follow you secretly to your mushroom patch, I can benefit from your knowledge and wisdom and skills, and maybe even catch that \textbf{fish} before you do.
Social \textbf{learning} really is visual theft.
And in any species that acquired it, it would behoove you to hide your best ideas, lest somebody steal them from you.
And so some time around 200, 000 years ago, our species confronted this crisis.
And we really had only two options for dealing with \textbf{the} conflicts that visual theft would bring.
One of those options was that we could have retreated into small family groups.
Because then \textbf{the} benefits of our ideas and knowledge would flow just to our \textbf{relatives}.
Had we chosen this option, sometime around 200, 000 years ago, we would probably still be living like \textbf{the} Neanderthals were when we first entered Europe 40, 000 years ago.
And this is because in small groups there are fewer ideas, there are fewer innovations.
And small groups are more prone to accidents and bad luck.
So if we ' d chosen that path, our \textbf{evolutionary} path would have led into \textbf{the} forest — and been a short one indeed. \textbf{The} other option we could choose was to develop \textbf{the} systems of communication that would allow us to share ideas and to cooperate amongst others.
Choosing this option would mean that a vastly greater fund of accumulated knowledge and wisdom would become available to any one individual than would ever arise from within an individual family or an individual person on their own.
Well, we chose \textbf{the} second option, and language is \textbf{the} result.
Language \textbf{evolved} to solve \textbf{the} crisis of visual theft.
Language is a piece of social technology for enhancing \textbf{the} benefits of cooperation — for reaching agreements, for striking deals and for coordinating our activities.
And you can see that, in a developing society that was beginning to acquire language, not having language would be a like a bird without \textbf{wings}.
Just as \textbf{wings} open up this \textbf{sphere} of air for birds to exploit, language opened up \textbf{the} \textbf{sphere} of cooperation for humans to exploit.
And we take this utterly for granted, because we ' re a species that is so at home with language, but you have to realize that even \textbf{the} simplest acts of exchange that we engage in are utterly dependent upon language.
And to see why, consider two scenarios from early in our \textbf{evolution}.
Let ' s imagine that you are really good at making arrowheads, but you ' re hopeless at making \textbf{the} wooden shafts with \textbf{the} flight feathers attached.
Two other people you know are very good at making \textbf{the} wooden shafts, but they ' re hopeless at making \textbf{the} arrowheads.
So what you do is — one of those people has not really acquired language yet.
And let ' s pretend \textbf{the} other one is good at language skills.
So what you do one day is you take a pile of arrowheads, and you walk up to \textbf{the} one that can ' t speak very well, and you put \textbf{the} arrowheads down in front of him, hoping that he ' ll get \textbf{the} idea that you want to trade your arrowheads for finished arrows.
But he looks at \textbf{the} pile of arrowheads, thinks they ' re a gift, picks them up, smiles and walks off.
Now you pursue this guy, gesticulating.
A scuffle ensues and you get stabbed with one of your own arrowheads. \textbf{Okay}, now replay this scene now, and you ' re approaching \textbf{the} one who has language.
You put down your arrowheads and say,"I ' d like to trade these arrowheads for finished arrows.
I ' ll split you 50 / 50."\textbf{The} other one says,"Fine.
Looks good to me.
We ' ll do that."Now \textbf{the} job is done.
Once we have language, we can put our ideas together and cooperate to have a prosperity that we couldn ' t have before we acquired it.
And this is why our species has prospered around \textbf{the} world while \textbf{the} rest of \textbf{the} animals sit behind bars in zoos, languishing.
That ' s why we build space shuttles and cathedrals while \textbf{the} rest of \textbf{the} world sticks sticks into \textbf{the} ground to extract termites.
All right, if this view of language and its value in solving \textbf{the} crisis of visual theft is true, any species that acquires it should show an \textbf{explosion} of creativity and prosperity.
And this is exactly what \textbf{the} archeological record shows.
If you look at our ancestors, \textbf{the} Neanderthals and \textbf{the} Homo erectus, our immediate ancestors, they ' re confined to small regions of \textbf{the} world.
But when our species arose about 200, 000 years ago, sometime after that we quickly walked out of Africa and spread around \textbf{the} entire world, occupying nearly every \textbf{habitat} on Earth.
Now whereas other species are confined to places that their genes adapt them to, with social \textbf{learning} and language, we could transform \textbf{the} environment to suit our needs.
And so we prospered in a way that no other animal has.
Language really is \textbf{the} most potent trait that has ever \textbf{evolved}.
It is \textbf{the} most valuable trait we have for converting new lands and resources into more people and their genes that natural \textbf{selection} has ever \textbf{devised}.
Language really is \textbf{the} voice of our genes.
Now having \textbf{evolved} language, though, we did something \textbf{peculiar}, even bizarre.
As we spread out around \textbf{the} world, we developed thousands of different languages.
Currently, there are about seven or 8, 000 different languages spoken on Earth.
Now you might say, well, this is just natural.
As we diverge, our languages are naturally going to diverge.
But \textbf{the} real puzzle and irony is that \textbf{the} greatest density of different languages on Earth is found where people are most tightly packed together.
If we go to \textbf{the} \textbf{island} of Papua New Guinea, we can find about 800 to 1, 000 distinct human languages, different human languages, spoken on that \textbf{island} alone.
There are places on that \textbf{island} where you can encounter a new language every two or three miles.
Now, incredible as this sounds, I once met a Papuan man, and I asked him if this could possibly be true.
And he said to me,"Oh no.
They ' re far closer together than that."And it ' s true ; there are places on that \textbf{island} where you can encounter a new language in under a mile.
And this is also true of some remote oceanic \textbf{islands}.
And so it seems that we use our language, not just to cooperate, but to draw rings around our cooperative groups and to establish identities, and perhaps to protect our knowledge and wisdom and skills from eavesdropping from outside.
And we know this because when we study different language groups and associate them with their cultures, we see that different languages slow \textbf{the} flow of ideas between groups.
They slow \textbf{the} flow of technologies.
And they even slow \textbf{the} flow of genes.
Now I can ' t speak for you, but it seems to be \textbf{the} case that we don ' t have sex with people we can ' t talk to.
Now we have to counter that, though, against \textbf{the} evidence we ' ve heard that we might have had some rather distasteful genetic dalliances with \textbf{the} Neanderthals and \textbf{the} Denisovans. \textbf{Okay}, this tendency we have, this seemingly natural tendency we have, towards isolation, towards keeping to ourselves, crashes head first into our modern world.
This remarkable image is not a map of \textbf{the} world.
In fact, it ' s a map of Facebook friendship links.
And when you plot those friendship links by their latitude and longitude, it literally draws a map of \textbf{the} world.
Our modern world is communicating with itself and with each other more than it has at any time in its past.
And that communication, that connectivity around \textbf{the} world, that globalization now raises a burden.
Because these different languages impose a barrier, as we ' ve just seen, to \textbf{the} transfer of goods and ideas and technologies and wisdom.
And they impose a barrier to cooperation.
And nowhere do we see that more clearly than in \textbf{the} European Union, whose 27 member countries speak 23 official languages. \textbf{The} European Union is now spending over one billion euros annually translating among their 23 official languages.
That ' s something on \textbf{the} order of 1. 45 billion U.
S. dollars on translation costs alone.
Now think of \textbf{the} absurdity of this situation.
If 27 individuals from those 27 member states sat around table, speaking their 23 languages, some very simple mathematics will tell you that you need an army of 253 translators to anticipate all \textbf{the} pairwise possibilities. \textbf{The} European Union employs a permanent staff of about 2, 500 translators.
And in 2007 alone — and I ' m sure there are more recent figures — something on \textbf{the} order of 1. 3 million pages were translated into English alone.
And so if language really is \textbf{the} solution to \textbf{the} crisis of visual theft, if language really is \textbf{the} conduit of our cooperation, \textbf{the} technology that our species derived to promote \textbf{the} free flow and exchange of ideas, in our modern world, we confront a question.
And that question is whether in this modern, globalized world we can really afford to have all these different languages.
To put it this way, nature knows no other circumstance in which functionally equivalent traits coexist.
One of them always drives \textbf{the} other extinct.
And we see this in \textbf{the} inexorable march towards standardization.
There are lots and lots of ways of measuring things — weighing them and measuring their length — but \textbf{the} metric system is winning.
There are lots and lots of ways of measuring time, but a really bizarre base 60 system known as hours and minutes and seconds is nearly universal around \textbf{the} world.
There are many, many ways of imprinting CDs or DVDs, but those are all being standardized as well.
And you can probably think of many, many more in your own everyday lives.
And so our modern world now is confronting us with a dilemma.
And it ' s \textbf{the} dilemma that this Chinese man faces, who ' s language is spoken by more people in \textbf{the} world than any other single language, and yet he is sitting at his blackboard, converting Chinese phrases into English language phrases.
And what this does is it raises \textbf{the} possibility to us that in a world in which we want to promote cooperation and exchange, and in a world that might be dependent more than ever before on cooperation to maintain and enhance our levels of prosperity, his actions suggest to us it might be inevitable that we have to confront \textbf{the} idea that our destiny is to be one world with one language.
Thank you.
Matt Ridley : Mark, one question.
Svante found that \textbf{the} FOXP2 gene, which seems to be associated with language, was also shared in \textbf{the} same form in Neanderthals as us.
Do we have any idea how we could have defeated Neanderthals if they also had language?
Mark Pagel : This is a very good question.
So many of you will be familiar with \textbf{the} idea that there ' s this gene called FOXP2 that seems to be implicated in some ways in \textbf{the} fine motor control that ' s associated with language. \textbf{The} reason why I don ' t believe that tells us that \textbf{the} Neanderthals had language is — here ' s a simple analogy : Ferraris are cars that have engines.
My car has an engine, but it ' s not a Ferrari.
Now \textbf{the} simple answer then is that genes alone don ' t, all by themselves, determine \textbf{the} outcome of very complicated things like language.
What we know about this FOXP2 and Neanderthals is that they may have had fine motor control of their mouths — who knows.
But that doesn ' t tell us they necessarily had language.
MR : Thank you very much indeed.

% matched lemmas: copy, device, devise, evolution, evolutionary, evolve, explosion, fish, fossil, habitat, island, learning, nut, object, okay, peculiar, pulse, relative, rock, selection, sphere, the, wing
\end{textsample}
